\documentclass{ximera}
%verbeterd 30/10/26
\title{Vraagstukken gemengde schakelingen: Basisoefeningen}
\author{Daan Lipkens}

\begin{document}
\begin{abstract}
    Basisoefeningen op de wet van Ohm, serieschakelingen en parallelschakelingen.
\end{abstract}
\maketitle

\section{Basisoefeningen schakelingen}

%———————————————————————————————————————
% BASIS 1: STROOMSTERKTE
\begin{exercise}
    Door een verbindingsdraad stroomt in $2{,}0\,\mathrm{min}$ een lading van $36\,\mathrm{C}$. Bereken de stroomsterkte.

    \[ I = \answer[tolerance=0.05,onlineshowanswerbutton]{0.30} \, \mathrm{A} \]

    \begin{solution}\nl
        \underline{Gegevens:}
        \begin{align*}
            & \Delta Q = 36\,\mathrm{C} & & \Delta t = 2{,}0\,\mathrm{min} = 120\,\mathrm{s}
        \end{align*}

        \underline{Gevraagd:} $I$ \\

        \underline{Oplossing:}
        \begin{align*}
            I &= \frac{\Delta Q}{\Delta t} = \frac{36\,\mathrm{C}}{120\,\mathrm{s}} = 0{,}30\,\mathrm{A}
        \end{align*}
    \end{solution}
\end{exercise}
%———————————————————————————————————————

%———————————————————————————————————————
% BASIS 2: OHM, STROOM
\begin{exercise}
    Een weerstand van $0{,}45\,\mathrm{k\Omega}$ wordt aangesloten op een spanning van $9{,}0\,\mathrm{V}$. Bereken de stroomsterkte in $\mathrm{mA}$.

    \[ I = \answer[tolerance=0.05,onlineshowanswerbutton]{20} \, \mathrm{mA} \]

    \begin{solution}\nl
        \underline{Gegevens:}
        \begin{align*}
            & R = 0{,}45\,\mathrm{k\Omega} = 450\,\mathrm{\Omega} & & U = 9{,}0\,\mathrm{V}
        \end{align*}

        \underline{Gevraagd:} $I$ \\

        \underline{Oplossing:}
        \begin{align*}
            I &= \frac{U}{R} = \frac{9{,}0\,\mathrm{V}}{450\,\mathrm{\Omega}} = 0{,}020\,\mathrm{A} = 20\,\mathrm{mA}
        \end{align*}
    \end{solution}
\end{exercise}
%———————————————————————————————————————

%———————————————————————————————————————
% BASIS 3: OHM, SPANNING
\begin{exercise}
    Door een weerstand van $48\,\mathrm{\Omega}$ loopt een stroom van $250\,\mathrm{mA}$. Bereken de spanning over de weerstand.

    \[ U = \answer[tolerance=0.05,onlineshowanswerbutton]{12} \, \mathrm{V} \]

    \begin{solution}\nl
        \underline{Gegevens:}
        \begin{align*}
            & R = 48\,\mathrm{\Omega} & & I = 250\,\mathrm{mA} = 0{,}250\,\mathrm{A}
        \end{align*}

        \underline{Gevraagd:} $U$ \\

        \underline{Oplossing:}
        \begin{align*}
            U &= R \cdot I = 48\,\mathrm{\Omega} \cdot 0{,}250\,\mathrm{A} = 12\,\mathrm{V}
        \end{align*}
    \end{solution}
\end{exercise}
%———————————————————————————————————————

%———————————————————————————————————————
% BASIS 4: OHM, WEERSTAND
\begin{exercise}
    Het verwarmingselement van een strijkijzer is aangesloten op $230\,\mathrm{V}$ en neemt een stroom van $0{,}46\,\mathrm{A}$ op. Bereken de weerstand van het element.

    \[ R = \answer[tolerance=0.05,onlineshowanswerbutton]{5.0} \cdot 10^{2} \, \mathrm{\Omega} \]

    \begin{solution}\nl
        \underline{Gegevens:}
        \begin{align*}
            & U = 230\,\mathrm{V} & & I = 0{,}46\,\mathrm{A}
        \end{align*}

        \underline{Gevraagd:} $R$ \\

        \underline{Oplossing:}
        \begin{align*}
            R &= \frac{U}{I} = \frac{230\,\mathrm{V}}{0{,}46\,\mathrm{A}} = 5{,}0 \cdot 10^{2}\,\mathrm{\Omega}
        \end{align*}
    \end{solution}
\end{exercise}
%———————————————————————————————————————

%———————————————————————————————————————
% BASIS 5: SERIE
\begin{exercise}
    Drie weerstanden $R_1 = 12\,\mathrm{\Omega}$, $R_2 = 18\,\mathrm{\Omega}$ en $R_3 = 30\,\mathrm{\Omega}$ staan in serie, aangesloten op een bron van $12\,\mathrm{V}$.

    \begin{center}
        \begin{circuitikz}[european]
            \draw (0,3) to[battery1, l=$U$] (0,0) -- (9,0) -- (9,3);
            \draw (0,3) to[R, l=$R_1$] (3,3) to[R, l=$R_2$] (6,3) to[R, l=$R_3$] (9,3);
        \end{circuitikz}
    \end{center}

    \begin{question}
        Bereken de vervangingsweerstand van de schakeling.

        \[ R_{\text{s}} = \answer[tolerance=0.05,onlineshowanswerbutton]{60} \, \mathrm{\Omega} \]

        \begin{solution}\nl
            \underline{Gegevens:}
            \begin{align*}
                & R_1 = 12\,\mathrm{\Omega} & & R_2 = 18\,\mathrm{\Omega} \\
                & R_3 = 30\,\mathrm{\Omega} & & U = 12\,\mathrm{V}
            \end{align*}

            \underline{Gevraagd:} $R_{\text{s}}$ \\

            \underline{Oplossing:}
            \begin{align*}
                R_{\text{s}} &= R_1 + R_2 + R_3 = 12\,\mathrm{\Omega} + 18\,\mathrm{\Omega} + 30\,\mathrm{\Omega} = 60\,\mathrm{\Omega}
            \end{align*}
        \end{solution}
    \end{question}

    \begin{question}
        Bereken de stroomsterkte door de bron.

        \[ I = \answer[tolerance=0.05,onlineshowanswerbutton]{0.20} \, \mathrm{A} \]

        \begin{solution}\nl
            \underline{Gevraagd:} $I$ \\

            \underline{Oplossing:}
            \begin{align*}
                I &= \frac{U}{R_{\text{s}}} = \frac{12\,\mathrm{V}}{60\,\mathrm{\Omega}} = 0{,}20\,\mathrm{A}
            \end{align*}
        \end{solution}
    \end{question}

    \begin{question}
        Bereken de spanning over $R_2$.

        \[ U_2 = \answer[tolerance=0.05,onlineshowanswerbutton]{3.6} \, \mathrm{V} \]

        \begin{hint}
            Wat weet je over de stroom door elke weerstand in een serieschakeling?
        \end{hint}

        \begin{solution}\nl
            \underline{Gevraagd:} $U_2$ \\

            \underline{Oplossing:}\\
            Door $R_2$ loopt dezelfde stroom als door de bron.
            \begin{align*}
                U_2 &= R_2 \cdot I = 18\,\mathrm{\Omega} \cdot 0{,}20\,\mathrm{A} = 3{,}6\,\mathrm{V}
            \end{align*}
        \end{solution}
    \end{question}
\end{exercise}
%———————————————————————————————————————

%———————————————————————————————————————
% BASIS 6: IDENTIEKE WEERSTANDEN PARALLEL
\begin{exercise}
    Vier identieke weerstanden van $1{,}2\,\mathrm{k\Omega}$ staan in parallel. Bereken de vervangingsweerstand.

    \[ R_{\text{p}} = \answer[tolerance=0.05,onlineshowanswerbutton]{3.0} \cdot 10^{2} \, \mathrm{\Omega} \]

    \begin{solution}\nl
        \underline{Gegevens:}
        \begin{align*}
            & R = 1{,}2\,\mathrm{k\Omega} = 1200\,\mathrm{\Omega} & & n = 4
        \end{align*}

        \underline{Gevraagd:} $R_{\text{p}}$ \\

        \underline{Oplossing:}
        \begin{align*}
            \frac{1}{R_{\text{p}}} &= 4 \cdot \frac{1}{1200\,\mathrm{\Omega}} = \frac{1}{300\,\mathrm{\Omega}} \\
            R_{\text{p}} &= 300\,\mathrm{\Omega} = 3{,}0 \cdot 10^{2}\,\mathrm{\Omega}
        \end{align*}
    \end{solution}
\end{exercise}
%———————————————————————————————————————

%———————————————————————————————————————
% BASIS 7: PARALLEL
\begin{exercise}
    Twee weerstanden $R_1 = 30\,\mathrm{\Omega}$ en $R_2 = 60\,\mathrm{\Omega}$ staan in parallel, aangesloten op een bron van $12\,\mathrm{V}$.

    \begin{center}
        \begin{circuitikz}[european]
            \draw (0,3) to[battery1, l=$U$] (0,0) -- (5,0);
            \draw (0,3) -- (5,3);
            \draw (2,3) to[R, l=$R_1$, *-*] (2,0);
            \draw (5,3) to[R, l=$R_2$, *-*] (5,0);
        \end{circuitikz}
    \end{center}

    \begin{question}
        Bereken de vervangingsweerstand van de schakeling.

        \[ R_{\text{p}} = \answer[tolerance=0.05,onlineshowanswerbutton]{20} \, \mathrm{\Omega} \]

        \begin{solution}\nl
            \underline{Gegevens:}
            \begin{align*}
                & R_1 = 30\,\mathrm{\Omega} & & R_2 = 60\,\mathrm{\Omega} \\
                & U = 12\,\mathrm{V} & &
            \end{align*}

            \underline{Gevraagd:} $R_{\text{p}}$ \\

            \underline{Oplossing:}
            \begin{align*}
                \frac{1}{R_{\text{p}}} &= \frac{1}{30\,\mathrm{\Omega}} + \frac{1}{60\,\mathrm{\Omega}} = \frac{2 + 1}{60\,\mathrm{\Omega}} = \frac{1}{20\,\mathrm{\Omega}} \\
                R_{\text{p}} &= 20\,\mathrm{\Omega}
            \end{align*}
        \end{solution}
    \end{question}

    \begin{question}
        Bereken de stroom door elke weerstand.

        \[ I_1 = \answer[tolerance=0.05,onlineshowanswerbutton]{0.40} \, \mathrm{A} \qquad I_2 = \answer[tolerance=0.05,onlineshowanswerbutton]{0.20} \, \mathrm{A} \]

        \begin{hint}
            Wat weet je over de spanning over elke tak in een parallelschakeling?
        \end{hint}

        \begin{solution}\nl
            \underline{Gevraagd:} $I_1$ en $I_2$ \\

            \underline{Oplossing:}\\
            Over beide weerstanden staat de bronspanning.
            \begin{align*}
                I_1 &= \frac{U}{R_1} = \frac{12\,\mathrm{V}}{30\,\mathrm{\Omega}} = 0{,}40\,\mathrm{A} \\
                I_2 &= \frac{U}{R_2} = \frac{12\,\mathrm{V}}{60\,\mathrm{\Omega}} = 0{,}20\,\mathrm{A}
            \end{align*}
        \end{solution}
    \end{question}

    \begin{question}
        Bereken de stroom door de bron.

        \[ I = \answer[tolerance=0.05,onlineshowanswerbutton]{0.60} \, \mathrm{A} \]

        \begin{solution}\nl
            \underline{Gevraagd:} $I$ \\

            \underline{Oplossing:}
            \begin{align*}
                I &= I_1 + I_2 = 0{,}40\,\mathrm{A} + 0{,}20\,\mathrm{A} = 0{,}60\,\mathrm{A}
            \end{align*}
            Controle: $\frac{U}{R_{\text{p}}} = \frac{12\,\mathrm{V}}{20\,\mathrm{\Omega}} = 0{,}60\,\mathrm{A}$.
        \end{solution}
    \end{question}
\end{exercise}
%———————————————————————————————————————

%———————————————————————————————————————
% BASIS 8: ONBEKENDE WEERSTAND PARALLEL
\begin{exercise}
    Een weerstand $R_1 = 30\,\mathrm{\Omega}$ staat in parallel met een onbekende weerstand $R_2$. De vervangingsweerstand van beide samen is $12\,\mathrm{\Omega}$. Bereken $R_2$.

    \[ R_2 = \answer[tolerance=0.05,onlineshowanswerbutton]{20} \, \mathrm{\Omega} \]

    \begin{hint}
        Schrijf eerst de vergelijking voor de vervangingsweerstand op met $R_2$ als onbekende. Welk deel kan je al invullen?
    \end{hint}

    \begin{solution}\nl
        \underline{Gegevens:}
        \begin{align*}
            & R_1 = 30\,\mathrm{\Omega} & & R_{\text{p}} = 12\,\mathrm{\Omega}
        \end{align*}

        \underline{Gevraagd:} $R_2$ \\

        \underline{Oplossing:}
        \begin{align*}
            \frac{1}{R_{\text{p}}} &= \frac{1}{R_1} + \frac{1}{R_2} \\
            \frac{1}{R_2} &= \frac{1}{R_{\text{p}}} - \frac{1}{R_1} = \frac{1}{12\,\mathrm{\Omega}} - \frac{1}{30\,\mathrm{\Omega}} = \frac{5 - 2}{60\,\mathrm{\Omega}} = \frac{1}{20\,\mathrm{\Omega}} \\
            R_2 &= 20\,\mathrm{\Omega}
        \end{align*}
    \end{solution}
\end{exercise}
%———————————————————————————————————————

\end{document}